\clearpage
\VerdierPhysicalPageAnchor{0173}%
\begingroup\small\setlength{\parskip}{0pt}\setlength{\emergencystretch}{1.5em}
\VerdierSemanticAnchor{ch4-4-4-4}
\medskip\noindent\textbf{4.4.4.}\quad
Soient $(H^n(p,q),\delta)$, $(H'{}^n(p,q),\delta')$ deux objets spectraux à
  valeurs dans $\mathcal A$ et
  $m^n(p,q)\colon\allowbreak H^n(p,q)\longrightarrow\allowbreak H'{}^n(p,q)$ un morphisme d'objets
spectraux. Un tel morphisme induit un morphisme des familles spectrales
correspondantes, i.e. une famille de morphismes :
\[
m\!\left[\begin{smallmatrix}n\\p,q,r,s\end{smallmatrix}\right]
\colon
E\!\left[\begin{smallmatrix}n\\p,q,r,s\end{smallmatrix}\right]
\longrightarrow
E'\!\left[\begin{smallmatrix}n\\p,q,r,s\end{smallmatrix}\right]
\]
qui commutent aux différentielles. En particulier, le morphisme $m(p,q)$
induit un morphisme des suites spectrales correspondantes :
\[
m_r^{p,q}\colon E_r^{p,q}\longrightarrow E'{}_{r}^{p,q}.
\]
(On pose ici $E_r^{p,q}=\mathrm I_r^{p,q}$ ou bien
$E_r^{p,q}=\mathrm{II}_r^{p,q}$ en prenant l'une ou l'autre des notations de
(4.3.3)).

\VerdierSemanticAnchor{ch4-4-4-5}
\medskip\noindent\textbf{Proposition 4.4.5.}\quad
\textit{On suppose que les objets spectraux}
\[
(H^n(p,q),\delta)
\qquad\textit{et}\qquad
(H'{}^n(p,q),\delta')
\]
\textit{sont stationnaires et que, pour un entier $r$ convenable, le
morphisme}
\[
m_r^{p,q}\colon E_r^{p,q}\longrightarrow E'{}_{r}^{p,q}
\]
\textit{est un isomorphisme pour tout couple $p,q$. Alors, pour tout
$r'\geq r$, le morphisme}
\[
m_{r'}^{p,q}\colon E_{r'}^{p,q}\longrightarrow E'{}_{r'}^{p,q}
\]
\textit{est un isomorphisme. Le morphisme}
\[
m^n(-\infty,+\infty)\colon H^n(-\infty,+\infty)
\longrightarrow H'{}^n(-\infty,+\infty)
\]
\textit{est un isomorphisme d'objets filtrés.}

\VerdierSemanticAnchor{ch4-4-4-5-proof}
En effet, les $m_r^{p,q}$ définissent un isomorphisme entre les objets
différentiels gradués $(E_r^{p,q},d_r)$ et $(E'{}_{r}^{p,q},d'_r)$. Ils
induisent donc un isomorphisme sur les objets de cohomologie
$E_{r+1}^{p,q}$ et $E'{}_{r+1}^{p,q}$, d'où la première assertion par
récurrence sur $r'$. Les objets spectraux étant stationnaires, le morphisme
\[
m_\infty^{p,q}\colon E_\infty^{p,q}\longrightarrow
E'{}_{\infty}^{p,q}
\]
est un isomorphisme. Les $E_\infty^{p,q}$ (resp. $E'{}_{\infty}^{p,q}$) sont
les gradués associés aux objets filtrés $H^n(-\infty,+\infty)$ (resp.
$H'{}^n(-\infty,+\infty)$). Comme les filtrations de ces objets sont
stationnaires, on en déduit que
\[
m^n(-\infty,+\infty)\colon H^n(-\infty,+\infty)
\longrightarrow H'{}^n(-\infty,+\infty)
\]
est un isomorphisme d'objets filtrés.
\endgroup
